% Springer Lecture Notes in Computer Science (LNCS) manuscript skeleton
% Project: Does Tokenization Fragmentation Predict Miscalibration in Bangla LLMs?
% Compile with the Springer LNCS class file (llncs.cls) and BibTeX.
% Keep this as the single main .tex manuscript; place figures beside this file.

\documentclass[runningheads]{llncs}

\usepackage[T1]{fontenc}
\usepackage{amsmath,amssymb}
\usepackage{graphicx}
\usepackage{booktabs}
\usepackage{multirow}
\usepackage{url}
\usepackage{hyperref}
\usepackage{xurl}
\hypersetup{
  colorlinks=true,
  linkcolor=black,
  citecolor=black,
  urlcolor=blue
}
\urlstyle{rm}

\begin{document}

\title{Does Tokenization Fragmentation Predict Miscalibration in Bangla LLMs? A Cross-Lingual Study}
\titlerunning{Tokenization Fragmentation and LLM Calibration}

% Author order follows the project README / team plan.
% TODO: Confirm spelling, author order, affiliations, email addresses, and ORCIDs with all authors.
\author{Ayetullah Khomane
  \and Md. Razaul Haque Subho
  \and Dipto Karmakar
  \and Naeem Hossain
  \and Abdullah al Azmi}
\authorrunning{Khomane et al.}

% TODO: Replace this placeholder with the correct institutional affiliations.
% Use \and between distinct affiliations and add \email{} where appropriate.
\institute{Affiliations and corresponding-author details to be confirmed.}

\maketitle

\begin{abstract}
% TODO: Write a 150--250 word abstract after the methodology and results are frozen.
% Include: motivation, research question, datasets/languages, model panel,
% fragmentation metrics, calibration metrics, main empirical findings, and contribution.
% Do not report preliminary branch results as final findings.
\end{abstract}

\keywords{Large language models \and Bangla \and Tokenization \and
Tokenization fragmentation \and Confidence calibration \and Cross-lingual evaluation}

\section{Introduction}
% Suggested content:
% 1. Background: LLM use in Bangla and other lower-resource languages.
% 2. Why tokenization fragmentation may matter for language representation.
% 3. Why confidence calibration matters in addition to accuracy.
% 4. Research problem, contributions, and paper structure.
% Keep the framing associational; do not imply causality.

\section{Background and Related Work}
\subsection{Tokenization and Fragmentation}
% Review tokenizer families, subword segmentation, language coverage,
% and prior fragmentation measures. Cite primary sources.

\subsection{Confidence Calibration in Large Language Models}
% Review confidence extraction, calibration evaluation, and limitations
% of ECE, Brier score, and other selected measures.

\subsection{Bangla Language Models and Cross-Lingual Evaluation}
% Cover relevant Bangla LLMs, BnMMLU, English comparison data,
% and the distinction between translated matched items and independent MMLU items.

\section{Research Questions}
% Confirm the final wording with the team before submission.
\begin{itemize}
  \item \textbf{RQ1:} How does tokenization fragmentation differ between Bangla and English for the selected LLM tokenizers?
  \item \textbf{RQ2:} Is tokenization fragmentation associated with confidence-calibration errors in Bangla LLM evaluations?
  \item \textbf{RQ3:} Does the observed association vary across models, subjects, or language conditions?
\end{itemize}

% TODO: Add hypotheses only if they are specified and justified in the study protocol.

\section{Models and Datasets}
\subsection{Datasets and Language Conditions}
% Document the exact dataset repository, configuration, split, version/revision,
% row filtering, duplicate handling, translation provenance, and item matching.
% Distinguish:
% (a) original Bangla BnMMLU items,
% (b) English translations of the same BnMMLU items,
% (c) independent English MMLU items, if used.
% Do not describe independent MMLU questions as matched translations.

\subsection{Model Selection and Identity}
% Insert the final frozen model panel. Record exact Hugging Face repository IDs,
% immutable revisions, tokenizer class, vocabulary size/hash, precision,
% context limit, and access date.
% The current branches contain different panels/model versions; reconcile before
% reporting a single definitive panel.

\subsection{Data Quality and Preprocessing}
% Explain schema validation, invalid rows, repeated questions, conflicting labels,
% Unicode normalization, whitespace handling, and any exclusions.
% Report counts from the exact dataset snapshot used in the final analysis.

\section{Tokenization Fragmentation Measurement}
\subsection{Tokenizer-Level Measurement}
% State whether the question alone or the full question-plus-options text is measured.
% Record tokenizer repository and revision. Do not load model weights for tokenizer-only
% audits unless the protocol explicitly requires model inference.

\subsection{Tokens per Word}
% Proposed primary measure; confirm the exact word-counting and normalization rules.
\begin{equation}
  \mathrm{TPW}(x) =
  \frac{N_{\mathrm{tokens}}(x)}
       {N_{\mathrm{words}}(x)}.
  \label{eq:tpw}
\end{equation}
% TODO: Define N_tokens, N_words, punctuation treatment, whitespace/token edge cases,
% special-token handling, empty strings, and Unicode normalization.

\subsection{Subword Token Retention Ratio}
% TODO: Include STRR only after its source definition and implementation are verified.
% State numerator, denominator, treatment of whole-word tokens, and interpretation.
% Do not invent or assume a definition from the acronym alone.

\subsection{Worked Tokenization Examples}
% Insert the team's checked Bangla-English examples.
% For each pair, report pair ID, exact text, tokenizer/model revision,
% token sequence or decoded pieces, token count, word count, and metric values.
% Keep model inference separate: tokenizer examples do not require inference.

\section{Confidence and Calibration Methodology}
\subsection{Answer Probability and Prediction}
% Specify how scores for options A--D are obtained, normalized, and converted
% into the predicted answer and confidence. Document prompt format, chat template,
% label-token handling, and any option-order rotation.

\subsection{Accuracy and Brier Score}
% Define multiclass accuracy and the exact Brier-score convention used.
\begin{equation}
  \mathrm{Brier} =
  \frac{1}{N}\sum_{i=1}^{N}\sum_{k=1}^{K}
  (p_{ik}-y_{ik})^2,
  \label{eq:brier}
\end{equation}
where $p_{ik}$ is the predicted probability for class $k$ on item $i$,
and $y_{ik}$ is the corresponding one-hot target.

\subsection{Expected Calibration Error}
% Freeze the number of bins, interval boundaries, endpoint convention,
% confidence definition, and empty-bin handling before final evaluation.
\begin{equation}
  \mathrm{ECE} =
  \sum_{b=1}^{B}\frac{|I_b|}{N}
  \left|\mathrm{acc}(I_b)-\mathrm{conf}(I_b)\right|.
  \label{eq:ece}
\end{equation}
% TODO: State whether raw or temperature-scaled probabilities are primary.
% If temperature scaling is used, document the fitting protocol and data separation.

\subsection{Reliability Analysis}
% Include reliability diagrams and any additional calibration diagnostics.
% Avoid asserting a universal ECE threshold for miscalibration.

\section{Experimental Setup}
\subsection{Evaluation Protocol}
% Record hardware, software/package versions, precision, batch size, random seed,
% prompt templates, checkpoint/resume behaviour, and model revision.
% Separate measured settings from estimates.

\subsection{Statistical Analysis}
% Specify the unit of analysis, model/subject controls, association tests or regression,
% confidence intervals, clustering, multiple-comparison correction, and missing-data rules.
% Account for repeated or duplicated question content where appropriate.
% Treat findings as observational associations, not causal effects.

\section{Results}
\subsection{Dataset and Model Coverage}
% TODO: Add final sample counts and model coverage table.

\subsection{Bangla-English Fragmentation Comparison}
% TODO: Add TPW/verified STRR summaries and paired comparisons.

\subsection{Accuracy and Calibration Results}
% TODO: Add accuracy, Brier, ECE, and reliability results from the frozen protocol.

\subsection{Fragmentation and Calibration Association}
% TODO: Report effect sizes, uncertainty intervals, and sensitivity analyses.
% Do not interpret descriptive trends as evidence of causation.

\section{Discussion}
% Interpret the findings against the research questions and related work.
% Discuss model-specific and language-specific patterns, alternative explanations,
% and the distinction between accuracy and calibration.

\section{Limitations}
% Discuss dataset translation quality, language coverage, model panel differences,
% prompt dependence, tokenizer comparability, sample dependence, and statistical power.
% Include any unresolved limitations supported by the final experiment.

\section{Conclusion and Future Work}
% Summarize what the study establishes, what remains uncertain,
% and the most useful next steps.

\begin{credits}
\subsubsection{\ackname}
% TODO: Add funding and acknowledgements, or remove this paragraph if not applicable.

\subsubsection{\discintname}
The authors declare that they have no competing interests relevant to this work.
% TODO: Confirm this statement with every author before submission.
\end{credits}

% Optional: Add a data/code availability statement if required by the venue.
% TODO: Provide stable repository release/commit and dataset access details.

\bibliographystyle{splncs04}
\bibliography{references}

\end{document}
